\section{Práctica de ingeniería y control}

\subsection{Planificación de Proof Search}
Hubert presentó un planificador: los Edges tienen distintos cost (Heuristic, Beam, MCTS), y en cada paso se elige la search policy en la tabla de chunk gating; chunk gating registra el token count, el Overlap y el proof term size. El planificador elige, según el latency budget, entre «search superficial + high replay» o «MCTS profundo + low replay».
\begin{table}[H]
\centering
\begin{tabular}{p{4cm}p{4cm}p{4cm}}
\toprule
Tipo de search & Características & Escenario de aplicación \\
\midrule
Heuristic tactics & Baja latency & exploration inicial \\
Beam search & selección conservadora de top-k & lemma discovery \\
MCTS + RL policy & value de estado + policy & polish final, adecuado para nivel de teorema \\
\bottomrule
\end{tabular}
\caption{Estrategia de planificación de Proof search}
\end{table}
\begin{importantbox}{KPI del planificador}
El latency budget, el proof-term length y el entropy drift determinan en conjunto el chunk gating gate; al superar el umbral, se degrada a conservative search.
\end{importantbox}

\subsection{Registro Trace y verificación}
Cada failed proof se registra en `trace.json`: incluye la tactic sequence, las indicaciones de Lean y el search depth. Hubert enfatizó que el trace debe tener un campo `human comment`, para que el governance team pueda hacer triage rápidamente. También señaló que, tras un fallo de verificación de Lean, la RL necesita replay de un alternate candidate para evitar quedar stuck.
\begin{warningbox}{Advertencia sobre el registro Trace}
Sin replay (trail) es imposible medir el strategy drift; un replay buffer sin trace hace que la policy entrene a ciegas con wrong proofs.
\end{warningbox}

\subsection{Gobernanza humano-máquina y reproducción}
El equipo de gobernanza envía periódicamente los high-drift chunk al human-in-the-loop board: el reviewer revisa el trace/log, evalúa la fairness y actualiza el knowledge summary en el newsletter. Esta práctica hace que AlphaProof trate cada fallo como seed de un experimento de reproducción.
\begin{knowledgebox}{Puntos clave de gobernanza}
1) High drift chunk activa un alert; 2) el reviewer genera un corrected tactic; 3) policy replay + reward correction; 4) el newsletter registra los puntos destacados de QA.
\end{knowledgebox}

\subsection{Resumen de la sección}
La parte de ingeniería destaca la planificación de Proof search, el registro de trace y el ciclo cerrado de gobernanza: mediante el triángulo de chunk gating, trace + review y newsletter, convierte al RL agent de una black box inexplicada en una línea vital de investigación y desarrollo reproducible.
\newpage

